% GENERATED FILE. Edit canonical lessons, then run npm run generate:books.
% canonical-source-hash: fnv1a64:c643de2794790ce8
% canonical-lessons: AR-C432-waqud, AR-C432-banzin, AR-C432-muharrik, AR-C432-faramil, AR-C432-zawwada

\chapter{Fuel and engines}
\label{ch:ar-fuel-and-engines}
\hlchaptermodality{432}

\begin{chapteropening}
\textbf{By the end of this chapter:} \emph{I can say fuel, petrol, an engine, brakes and to supply.}

Complete the last lesson of chapter 432: I can say fuel, petrol, an engine, brakes and to supply.
\end{chapteropening}

\section[\texorpdfstring{\ar{وقود} (waqūd)}{وقود (waqūd)}]{\ar{وقود} (waqūd) — fuel}

\label{lesson:AR-C432-waqud}



\begin{quote}
Before the new one: say the Arabic for congestion, then the Arabic for to turn off.
\end{quote}

\subsection*{You'll want to know: \ar{وقود}}
\textbf{\ar{وقود}} — \emph{waqūd} — \textquotedblleft{}fuel\textquotedblright{}.

\begin{quote}
\ar{اشترينا} \ar{الوقود} \ar{من} \ar{محطة} \ar{صغيرة}.
\end{quote}

\begin{quote}
\emph{ishtaraynā al-waqūd min maḥaṭṭa ṣaghīra.}
\end{quote}

\begin{quote}
We bought fuel at a small station.
\end{quote}

\begin{cousinweb}
\textbf{\ar{وقود}} (\emph{waqūd}) is built on the native Arabic root \emph{w-q-d}, \textquotedblleft{}to burn, to blaze\textquotedblright{} (Lane). A \emph{waqūd} is what is burned.
\end{cousinweb}

\begin{grammarlens}[title={What you've built}]
One more word for getting around.
\end{grammarlens}

\subsection*{Your turn}
\begin{itemize}
\raggedright
  \item \emph{Say it:} \emph{waqūd}
  \item \emph{Say it:} \emph{waqūd}, once more
  \item \emph{Say it:} \emph{ishtaraynā al-waqūd min maḥaṭṭa ṣaghīra.} — We bought fuel at a small station.
  \item \emph{Recall:} say the Arabic for congestion, then the Arabic for to turn off, then say \emph{waqūd} again
\end{itemize}

\begin{tcolorbox}[breakable,colback=teal!4,colframe=teal!35!black,title={Before you move on}]
What does \ar{وقود} mean? (\textquotedblleft{}Fuel\textquotedblright{}.) Say it once more.
\end{tcolorbox}
\section[\texorpdfstring{\ar{بنزين} (banzīn)}{بنزين (banzīn)}]{\ar{بنزين} (banzīn) — petrol}

\label{lesson:AR-C432-banzin}



\begin{quote}
Before the new one: say the Arabic for to turn off, then the Arabic for fuel.
\end{quote}

\subsection*{You'll want to know: \ar{بنزين}}
\textbf{\ar{بنزين}} — \emph{banzīn} — \textquotedblleft{}petrol\textquotedblright{}.

\begin{quote}
\ar{كانت} \ar{السيارة} \ar{بحاجة} \ar{إلى} \ar{البنزين}.
\end{quote}

\begin{quote}
\emph{kānat as-sayyāra bi-ḥāja ilā al-banzīn.}
\end{quote}

\begin{quote}
The car needed petrol.
\end{quote}

\begin{cousinweb}
\textbf{\ar{بنزين}} (\emph{banzīn}) was borrowed from German \emph{Benzin} (EtymDB).
\end{cousinweb}

\begin{grammarlens}[title={What you've built}]
One more word for getting around.
\end{grammarlens}

\subsection*{Your turn}
\begin{itemize}
\raggedright
  \item \emph{Say it:} \emph{banzīn}
  \item \emph{Say it:} \emph{banzīn}, once more
  \item \emph{Say it:} \emph{kānat as-sayyāra bi-ḥāja ilā al-banzīn.} — The car needed petrol.
  \item \emph{Recall:} say the Arabic for to turn off, then the Arabic for fuel, then say \emph{banzīn} again
\end{itemize}

\begin{tcolorbox}[breakable,colback=teal!4,colframe=teal!35!black,title={Before you move on}]
What does \ar{بنزين} mean? (\textquotedblleft{}Petrol\textquotedblright{}.) Say it once more.
\end{tcolorbox}
\section[\texorpdfstring{\ar{محرك} (muḥarrik)}{محرك (muḥarrik)}]{\ar{محرك} (muḥarrik) — an engine; a motor}

\label{lesson:AR-C432-muharrik}



\begin{quote}
Before the new one: say the Arabic for fuel, then the Arabic for petrol.
\end{quote}

\subsection*{You'll want to know: \ar{محرك}}
\textbf{\ar{محرك}} — \emph{muḥarrik} — \textquotedblleft{}an engine; a motor\textquotedblright{}.

\begin{quote}
\ar{كان} \ar{محرك} \ar{السيارة} \ar{صاخبا}.
\end{quote}

\begin{quote}
\emph{kāna muḥarrik as-sayyāra ṣākhiban.}
\end{quote}

\begin{quote}
The car's engine was noisy.
\end{quote}

\begin{cousinweb}
\textbf{\ar{محرك}} (\emph{muḥarrik}) is built on the native Arabic root \emph{ḥ-r-k}, \textquotedblleft{}to move\textquotedblright{} (Lane). It is the Form II participle, \textquotedblleft{}the thing that sets something moving\textquotedblright{}.
\end{cousinweb}

\begin{grammarlens}[title={What you've built}]
One more word for getting around.
\end{grammarlens}

\subsection*{Your turn}
\begin{itemize}
\raggedright
  \item \emph{Say it:} \emph{muḥarrik}
  \item \emph{Say it:} \emph{muḥarrik}, once more
  \item \emph{Say it:} \emph{kāna muḥarrik as-sayyāra ṣākhiban.} — The car's engine was noisy.
  \item \emph{Recall:} say the Arabic for fuel, then the Arabic for petrol, then say \emph{muḥarrik} again
\end{itemize}

\begin{tcolorbox}[breakable,colback=teal!4,colframe=teal!35!black,title={Before you move on}]
What does \ar{محرك} mean? (\textquotedblleft{}An engine; a motor\textquotedblright{}.) Say it once more.
\end{tcolorbox}
\section[\texorpdfstring{\ar{فرامل} (farāmil)}{فرامل (farāmil)}]{\ar{فرامل} (farāmil) — brakes}

\label{lesson:AR-C432-faramil}



\begin{quote}
Before the new one: say the Arabic for petrol, then the Arabic for an engine.
\end{quote}

\subsection*{You'll want to know: \ar{فرامل}}
\textbf{\ar{فرامل}} — \emph{farāmil} — \textquotedblleft{}brakes\textquotedblright{}.

\begin{quote}
\ar{ضغط} \ar{أبي} \ar{على} \ar{الفرامل} \ar{بسرعة}.
\end{quote}

\begin{quote}
\emph{ḍaghaṭa abī ʿalā al-farāmil bi-surʿa.}
\end{quote}

\begin{quote}
My father pressed the brakes quickly.
\end{quote}

\begin{cousinweb}
\textbf{\ar{فرامل}} (\emph{farāmil}) is a loanword whose origin is not settled; it is often traced to Italian \emph{freno}, \textquotedblleft{}a brake\textquotedblright{}.
\end{cousinweb}

\begin{grammarlens}[title={What you've built}]
One more word for getting around.
\end{grammarlens}

\subsection*{Your turn}
\begin{itemize}
\raggedright
  \item \emph{Say it:} \emph{farāmil}
  \item \emph{Say it:} \emph{farāmil}, once more
  \item \emph{Say it:} \emph{ḍaghaṭa abī ʿalā al-farāmil bi-surʿa.} — My father pressed the brakes quickly.
  \item \emph{Recall:} say the Arabic for petrol, then the Arabic for an engine, then say \emph{farāmil} again
\end{itemize}

\begin{tcolorbox}[breakable,colback=teal!4,colframe=teal!35!black,title={Before you move on}]
What does \ar{فرامل} mean? (\textquotedblleft{}Brakes\textquotedblright{}.) Say it once more.
\end{tcolorbox}
\section[\texorpdfstring{\ar{زود} (zawwada)}{زود (zawwada)}]{\ar{زود} (zawwada) — to supply; to fill up (with)}

\label{lesson:AR-C432-zawwada}



\begin{quote}
Before the new one: say the Arabic for an engine, then the Arabic for brakes.
\end{quote}

\subsection*{You'll want to know: \ar{زود}}
\textbf{\ar{زود}} — \emph{zawwada} — \textquotedblleft{}to supply; to fill up (with)\textquotedblright{}.

\begin{quote}
\ar{زود} \ar{العامل} \ar{السيارة} \ar{بالبنزين}.
\end{quote}

\begin{quote}
\emph{zawwada al-ʿāmil as-sayyāra bi-l-banzīn.}
\end{quote}

\begin{quote}
The worker filled the car up with petrol.
\end{quote}

Present: \ar{يزود}, \emph{yuzawwid}.

It takes \emph{bi-} before the thing supplied.

\begin{cousinweb}
\textbf{\ar{زود}} (\emph{zawwada}) is built on the native Arabic root \emph{z-w-d}, \textquotedblleft{}provisions for a journey\textquotedblright{} (Lane). It is Form II, \textquotedblleft{}to give someone provisions\textquotedblright{}: the root of \emph{zād}, food for the road.
\end{cousinweb}

\begin{grammarlens}[title={What you've built}]
One more word for getting around.
\end{grammarlens}

\subsection*{Your turn}
\begin{itemize}
\raggedright
  \item \emph{Say it:} \emph{zawwada}
  \item \emph{Say it:} \emph{zawwada}, once more
  \item \emph{Say it:} \emph{zawwada al-ʿāmil as-sayyāra bi-l-banzīn.} — The worker filled the car up with petrol.
  \item \emph{Recall:} say the Arabic for an engine, then the Arabic for brakes, then say \emph{zawwada} again
\end{itemize}

\begin{tcolorbox}[breakable,colback=teal!4,colframe=teal!35!black,title={Before you move on}]
What does \ar{زود} mean? (\textquotedblleft{}To supply; to fill up (with)\textquotedblright{}.) Say it once more.
\end{tcolorbox}
